\section{Feature Requests (geplante Erweiterungen)}
\label{sec:featurerequests}

Waehrend User Stories den \textit{aktuellen} Funktionsumfang beschreiben, dokumentieren
Feature Requests die \textit{zukuenftigen} Erweiterungsmoeglichkeiten. Sie entstehen aus
eigenen Ideen, Rueckmeldungen potenzieller Nutzer und technischen Ueberlegungen.

Die Priorisierung erfolgt nach dem MoSCoW-Schema:
\begin{itemize}
  \item \textbf{Must} -- Wird in der naechsten Version umgesetzt
  \item \textbf{Should} -- Sollte bald umgesetzt werden
  \item \textbf{Could} -- Waere schoen, aber nicht dringend
  \item \textbf{Won't (yet)} -- Aktuell nicht geplant, aber vorstellbar
\end{itemize}

\subsection{FR-01: Automatischer PDF-Versand per E-Mail}
\begin{description}
  \item[Prioritaet:] Should
  \item[Beschreibung:] Der Benutzer soll einen Zeitplan konfigurieren koennen (z.\,B.\
    woechtentlich oder monatlich), nach dem SignalReport automatisch einen PDF-Bericht
    generiert und per SMTP an eine konfigurierte E-Mail-Adresse versendet.
  \item[Motivation:] Viele Benutzer moechten regelmaessige Berichte erhalten, ohne sich
    manuell einloggen zu muessen. Besonders fuer Provider-Beschwerden ist ein monatlicher
    Report hilfreich.
  \item[Technische Ueberlegungen:] SMTP-Konfiguration in config.json; javax.mail oder
    Jakarta Mail; Scheduler via ScheduledExecutorService.
\end{description}

\subsection{FR-02: Prometheus-Exporter fuer Grafana-Integration}
\begin{description}
  \item[Prioritaet:] Could
  \item[Beschreibung:] Ein \texttt{/metrics}-Endpoint im Prometheus-Exposition-Format soll
    bereitgestellt werden, damit bestehende Monitoring-Infrastrukturen (Grafana, Prometheus)
    die SignalReport-Daten integrieren koennen.
  \item[Motivation:] Fortgeschrittene Benutzer haben oft bereits ein Grafana-Dashboard und
    moechten die Internetqualitaet dort einbinden.
  \item[Technische Ueberlegungen:] Einfaches Text-Format; keine externe Library noetig;
    Metriken: \texttt{signalreport\_ping\_latency\_ms}, \texttt{signalreport\_dns\_latency\_ms},
    \texttt{signalreport\_http\_latency\_ms}, \texttt{signalreport\_packet\_loss\_percent}.
\end{description}

\subsection{FR-03: Multi-Host-Vergleich}
\begin{description}
  \item[Prioritaet:] Could
  \item[Beschreibung:] Wenn SignalReport auf mehreren Geraeten laeuft (z.\,B.\ Desktop und NAS),
    sollen die Messdaten in einer zentralen Ansicht verglichen werden koennen.
  \item[Motivation:] Ermoeglicht die Unterscheidung, ob ein Problem am Geraet (WLAN vs.\
    Kabel) oder am Provider liegt.
  \item[Technische Ueberlegungen:] Bereits vorbereitet durch das Host-Hash-System in der
    Datenbank; erfordert API-Endpunkte fuer Host-uebergreifende Abfragen.
\end{description}

\subsection{FR-04: Docker-Container}
\begin{description}
  \item[Prioritaet:] Should
  \item[Beschreibung:] Ein offizielles Docker-Image soll bereitgestellt werden, das
    SignalReport als Container ausfuehrbar macht (z.\,B.\ auf einem NAS mit Docker-Support
    oder in einer Cloud-Umgebung).
  \item[Motivation:] Docker ist der De-facto-Standard fuer einfache Deployments; besonders
    auf Synology NAS mit Container Manager.
  \item[Technische Ueberlegungen:] Dockerfile basierend auf \texttt{eclipse-temurin:21-jre};
    Volume-Mount fuer \texttt{data/} und \texttt{config.json}; Port-Mapping 4567.
\end{description}

\subsection{FR-05: Traceroute-Analyse}
\begin{description}
  \item[Prioritaet:] Won't (yet)
  \item[Beschreibung:] Zusaetzlich zur reinen Latenz-Messung soll ein Traceroute zum Ziel
    durchgefuehrt und visualisiert werden koennen, um den genauen Punkt im Netzwerkpfad zu
    identifizieren, an dem Verzoegerungen entstehen.
  \item[Motivation:] Hilft bei der Diagnose, ob das Problem beim eigenen Router, beim
    Provider oder auf dem Weg zum Zielserver liegt.
  \item[Technische Ueberlegungen:] Plattformabhaengig (traceroute/tracert); erfordert
    erhoehte Rechte auf manchen Systemen; lange Laufzeit.
\end{description}

\subsection{FR-06: Speedtest-Integration}
\begin{description}
  \item[Prioritaet:] Won't (yet)
  \item[Beschreibung:] Neben Latenz-Messungen soll auch die tatsaechliche Bandbreite
    (Download/Upload) gemessen werden koennen.
  \item[Motivation:] Fuer Provider-Beschwerden ist die Bandbreite oft relevanter als die
    Latenz. Viele Nutzer erwarten einen \enquote{Speedtest} als Grundfunktion.
  \item[Technische Ueberlegungen:] Erfordert definierte Test-Server oder eine API wie
    Ookla/Cloudflare Speedtest; erzeugt signifikanten Traffic; nicht fuer kurze Intervalle
    geeignet.
\end{description}

\subsection{Zusammenfassung der Feature Requests}

\begin{table}[h]
\centering
\begin{tabular}{|l|l|l|}
\hline
\textbf{ID} & \textbf{Feature} & \textbf{Prioritaet} \\
\hline
FR-01 & Automatischer PDF-Versand per E-Mail & Should \\
FR-02 & Prometheus-Exporter (Grafana) & Could \\
FR-03 & Multi-Host-Vergleich & Could \\
FR-04 & Docker-Container & Should \\
FR-05 & Traceroute-Analyse & Won't (yet) \\
FR-06 & Speedtest-Integration & Won't (yet) \\
\hline
\end{tabular}
\caption{Feature Requests mit MoSCoW-Priorisierung}
\label{tab:featurerequests}
\end{table}
